\scanpage{111}
OPS $f\deleted{\unsure{p}}$ de la forme $p'+\alpha p+\beta$ (quitte à multiplier par un scalaire)
\[
 \left\{\begin{aligned}
 f&=p'+\alpha p+\beta,\quad \alpha,\beta\in k\\
 \operatorname{div} f&=\{x\}+\{y\}+\{z\}-3\{0_E\}
 \end{aligned}\right. \tag*{(43)}
\]
On aura donc
\[
 f(x)=f(y)=f(z)=0\qquad\text{i.e.} \tag*{(44)}
\]
\[
 \left\{\begin{aligned}
 p'(x)+\alpha p(x)+\beta&=0\\
 p'\unsure{y)|(y)}+\alpha p(y)+\beta&=0\\
 p'(z)+\alpha p(z)+\beta&=0
 \end{aligned}\right. \tag*{(45)}
\]
ce qui signifie aussi que le vecteur colonne $\begin{pmatrix}p'(x)\\p'(y)\\p'(z)\end{pmatrix}$ s'exprime comme comb. lin. des vecteurs colonnes $\begin{pmatrix}p(x)\\p(y)\\p(z)\end{pmatrix}$ et $\begin{pmatrix}1\\1\\1\end{pmatrix}$, \inserted{ce} qui \deleted{\unsure{don}}implique que ces trois vect. colonnes sont lin. dépendants, d'où (42). Inversement, supposons qu'on ait (42), et $x,y,z$ \deleted{\unsure{deux|égaux}} \inserted{tous distincts}, \deleted{\unsure{par ex|pas}} \deleted{que l'on} \inserted{alors on} n'a pas $p(x)=p(y)=p(z)$, \deleted{\unsure{sinon}} \inserted{(sinon $x,y,z$ appartiendraient à une $=$ fibre de $p$, donc deux parmi eux seraient égaux)}, \deleted{\unsure{et que}} \deleted{\unsure{alors}} \inserted{donc} les vect. colonne $\begin{pmatrix}1\\1\\1\end{pmatrix}$ $\begin{pmatrix}p(x)\\p(y)\\p(z)\end{pmatrix}$ sont lin. indép., et (42) signifie qu'il existe $\alpha,\beta\in k$ tels qu'on ait (45), i.e. (44), où $f$ est donnée par (43) a). Comme $x,y,z$ sont distincts, on en conclut bien (43) b), d'où (41).
% archive sheet number: 110
